% !TEX TS-program = pdflatex
% !TeX program = pdflatex
% !TEX encoding = UTF-8
% !TEX spellcheck = fr

\documentclass[xcolor=table]{beamer}


%\usepackage{fullpage}
%\usepackage[left=2.8cm,right=2.2cm,top=2 cm,bottom=2 cm]{geometry}
\setbeamersize{text margin left=10pt,text margin right=10pt}


\usepackage{amsmath,amssymb} 
\usepackage{fontenc}
\usepackage{textcomp}
\usepackage[utf8]{inputenc}
\usepackage{CJKutf8}
\usepackage[french]{babel}
\usepackage{arabtex,acolor}
\usepackage{txfonts}
\usepackage{tipa}
\usepackage[]{graphicx}
\usepackage{multirow}
\usepackage{hyperref}
\usepackage{colortbl}
\usepackage{tabularray}
\usepackage{listingsutf8}
\usepackage{wrapfig}
\usepackage{multicol}
\usepackage[export]{adjustbox} %for images in table, also for frame
\usepackage[many]{tcolorbox}
\usepackage{wasysym}
\usepackage[french,lined]{algorithm2e}
\usepackage{alltt} %verbatim with commands
\usepackage{longtable}
\usepackage{tabu}


\definecolor{lightblue}{HTML}{D0D2FF}
\definecolor{lightyellow}{HTML}{FFFFAA}
\definecolor{darkblue}{HTML}{0000BB}
\definecolor{olivegreen}{HTML}{006600}
\definecolor{darkgreen}{HTML}{008B45} %009B55
\definecolor{violet}{HTML}{6600CC}
%\definecolor{deeppink}{HTML}{FF1493}
\definecolor{orangey}{HTML}{FFBB00}

\usetheme{Warsaw} % Antibes Boadilla Warsaw

\beamertemplatenavigationsymbolsempty


\let\oldcite\cite
\renewcommand{\cite}[1]{{\bfseries\color{orangey}\oldcite{#1}}}

\tcbuselibrary{listings}

\newcommand{\kurl}[1]{{\scriptsize\bfseries\color{orangey}\url{#1}}}


%\renewcommand{\baselinestretch}{1.5}

\def\supit#1{\raisebox{0.8ex}{\small\it #1}\hspace{0.05em}}

\AtBeginSection{%
	\begin{frame}
		\sectionpage
	\end{frame}
}

\newcommand{\rottext}[2]{%
	\rotatebox{90}{%
	\begin{minipage}{#1}%
		\raggedleft#2%
	\end{minipage}%
	}%
}


\institute{ %
Laboratoire de la Communication dans les Systèmes Informatiques (LCSI)

\vspace{6pt}
École  nationale Supérieure d'Informatique (ESI, ex. INI), Alger, Algérie
}
\author[ \textbf{\footnotesize\insertframenumber/\inserttotalframenumber} \hspace*{1cm} ESI - ARIES Abdelkrime (Master 2022/2023)] %
{ARIES Abdelkrime}
%\titlegraphic{\includegraphics[height=1cm]{../img/esi-logo.png}%\hspace*{4.75cm}~
	
	
	\date{Master 2022/2023} %\today

\titlegraphic{%
	\includegraphics[height=1cm]{../img/esi-logo.png}%
	\hspace{2cm}
	\includegraphics[height=1cm]{../img/lcsi-logo.png}%
}


%\setbeamertemplate{headline}{}



\newcommand{\keyword}[1]{\textcolor{red}{\bfseries\itshape #1}}
\newcommand{\expword}[1]{\textcolor{olivegreen}{#1}}
\newcommand{\optword}[1]{\textcolor{violet}{\bfseries #1}}

\makeatletter
\newcommand\mysphere{%
	\parbox[t]{10pt}{\raisebox{0.2pt}{\beamer@usesphere{item projected}{bigsphere}}}}
\makeatother

%\let\oldtabular\tabular
%\let\endoldtabular\endtabular
%\renewenvironment{tabular}{\rowcolors{2}{white}{lightblue}\oldtabular\rowcolor{blue}}{\endoldtabular}


\NoAutoSpacing %french autospacing after ":"

\def\graphpath{}

\newcommand{\changegraphpath}[1]{\def\graphpath{#1}}


\newcommand{\vgraphpage}[2][.84\textheight]{%
%	\begin{center}%
		\includegraphics[height=#1]{\graphpath #2}%
%	\end{center}%
}

\newcommand{\hgraphpage}[2][\textwidth]{%
%	\begin{center}%
		\includegraphics[width=#1]{\graphpath #2}%
%	\end{center}%
}

\newcommand{\graphpage}[2][]{%
	\includegraphics[#1]{\graphpath #2}%
}

\bibliographystyle{apalike}

\newcommand{\insertbibliography}[2]{
	\appendix
	\section*{Bibliographie}
	\nocite{#2}
%	\makeatletter % to change template
%	\setbeamertemplate{headline}[default] % not mandatory, but I though it was better to set it blank
%	\def\beamer@entrycode{\vspace*{-\headheight}} % here is the part we are interested in :)
%	\makeatother
	\begin{multicols*}{2}[\frametitle{\insertsection} \usebeamertemplate*{frametitle}]%\usebeamertemplate*{frametitle}\frametitle{Références}
		\tiny
		\bibliography{#1}
	\end{multicols*}
}

\definecolor{my-grey}{RGB}{233, 233, 233}

\newcommand{\insertlicence}{
	\begin{frame}[plain]
	\frametitle{Licence : CC-BY 4.0}
%	\framesubtitle{Licence: CC-BY-NC 4.0}

	\begin{tcolorbox}[colback=cyan,
		colframe=cyan,  
		arc=0pt,outer arc=0pt,
		valign=top, 
		halign=center,
		width=\textwidth]
		
		\includegraphics[width=.5cm]{../img/licence/cc_icon_white_x2.png}
		\includegraphics[width=.5cm]{../img/licence/attribution_icon_white_x2.png}
		
		\color{white}
		\bfseries Attribution 4.0 International (CC BY 4.0) \\
		\tiny \url{https://creativecommons.org/licenses/by/4.0/deed.fr}
		
	\end{tcolorbox}\vspace{-.5cm}
	\begin{tcolorbox}[colback=my-grey,
		colframe=my-grey,  
		center, arc=0pt,outer arc=0pt,
		valign=top, 
		halign=left,
		width=\textwidth]
		
		\tiny
		
		\begin{center}
			\bfseries\Large
			Vous êtes autorisé à :
		\end{center}
		
		\begin{minipage}{0.83\textwidth}
			\begin{itemize}
				\item[] \textbf{Partager} — copier, distribuer et communiquer le matériel par tous moyens et sous tous formats
				\item[] \textbf{Adapter} — remixer, transformer et créer à partir du matériel
				pour toute utilisation, y compris commerciale.
			\end{itemize}
		\end{minipage}
		\begin{minipage}{0.15\textwidth}
			\includegraphics[width=\textwidth]{../img/licence/FreeCulturalWorks_seal_x2.jpg}
		\end{minipage}
	
		
		\begin{center}
			\bfseries\Large
			Selon les conditions suivantes :
		\end{center}
		
		\begin{itemize}
			\item[] \textbf{Attribution} — Vous devez créditer l'Œuvre, intégrer un lien vers la licence et indiquer si des modifications ont été effectuées à l'Oeuvre. Vous devez indiquer ces informations par tous les moyens raisonnables, sans toutefois suggérer que l'Offrant vous soutient ou soutient la façon dont vous avez utilisé son Oeuvre. 
			\item[] \textbf{Pas de restrictions complémentaires} — Vous n'êtes pas autorisé à appliquer des conditions légales ou des mesures techniques qui restreindraient légalement autrui à utiliser l'Oeuvre dans les conditions décrites par la licence.
		\end{itemize}
		
	\end{tcolorbox}
	
%	\begin{center}
%		\bfseries Attribution 4.0 International (CC BY 4.0)
%		\url{https://creativecommons.org/licenses/by/4.0/deed.fr}
%	\end{center}

%	\tiny
%
%	Vous êtes autorisé à : 
%	\begin{itemize}
%		\item \textbf{Partager} — copier, distribuer et communiquer le matériel par tous moyens et sous tous formats
%		\item \textbf{Adapter} — remixer, transformer et créer à partir du matériel
%	\end{itemize}
%	
%	Selon les conditions suivantes : 
%	\begin{itemize}
%		\item \textbf{Attribution} — Vous devez créditer l'Œuvre, intégrer un lien vers la licence et indiquer si des modifications ont été effectuées à l'Oeuvre. Vous devez indiquer ces informations par tous les moyens raisonnables, sans toutefois suggérer que l'Offrant vous soutient ou soutient la façon dont vous avez utilisé son Oeuvre.
%		\item \textbf{Pas d'Utilisation Commerciale} — Vous n'êtes pas autorisé à faire un usage commercial de cette Oeuvre, tout ou partie du matériel la composant. 
%		\item \textbf{Pas de restrictions complémentaires} — Vous n'êtes pas autorisé à appliquer des conditions légales ou des mesures techniques qui restreindraient légalement autrui à utiliser l'Oeuvre dans les conditions décrites par la licence.
%	\end{itemize}

	\end{frame}
}

\settowidth{\leftmargini}{\usebeamertemplate{itemize item}}
\addtolength{\leftmargini}{\labelsep}

\AtBeginDocument{
	\newcolumntype{L}[2]{>{\vbox to #2\bgroup\vfill\flushleft}p{#1}<{\egroup}} 
	
	\begin{frame}[plain]
		\maketitle
	\end{frame}

	\insertlicence
}


% needs etoolbox; to break links after -
\appto\UrlBreaks{\do\-}


\makeatletter
\newcommand{\xRightarrow}[2][]{\ext@arrow 0359\Rightarrowfill@{#1}{#2}}
\makeatother


\usefonttheme{structurebold}
%\usefonttheme{professionalfonts}

\title[TALN : 06- Sens et Désambigüisation]%
{Traitement automatique du langage naturel\\Chapitre 06 : Sens des mots et désambigüisation lexicale} 

\changegraphpath{../img/sens/}

\begin{document}
	
\begin{frame}
\frametitle{Traitement automatique du langage naturel}
\framesubtitle{Sens des mots et désambigüisation lexicale : Introduction}

%\begin{exampleblock}{Exemple d'une phrase en français}
%	\begin{center}
%		\Large\bfseries
%		L'élève a écrit une solution avec un stylo. 
%		
%		L'élève a écrit une solution avec une explication.
%	\end{center}
%\end{exampleblock}

\begin{itemize}
	\item 
\end{itemize}

\end{frame}

%\begin{frame}
%\frametitle{Traitement automatique du langage naturel}
%\framesubtitle{Sens des mots et désambigüisation lexicale : Un peu d'humour}
%
%\begin{center}
%	\vgraphpage{humour-parse.jpg}
%\end{center}
%
%\end{frame}

\begin{frame}
\frametitle{Traitement automatique du langage naturel}
\framesubtitle{Sens des mots et désambigüisation lexicale : Plan}

\begin{multicols}{2}
%	\small
\tableofcontents
\end{multicols}
\end{frame}

%===================================================================================
\section{Représentation vectorielle des mots}
%===================================================================================

\begin{frame}
\frametitle{Sens des mots et désambigüisation lexicale}
\framesubtitle{Représentation vectorielle des mots}

\begin{itemize}
	\item 	
\end{itemize}

\end{frame}


\subsection{Tf-Idf}

\begin{frame}
\frametitle{Sens et désambigüisation : Représentation vectorielle}
\framesubtitle{Tf-idf}

\begin{itemize}
	\item Le sens d'un document/phrase peut être représenté par ces mots 
	\item La fréquence d'un mot dans un document/phrase s'appelle \keywords{term frequency (TF)}
	\item Il y a des mots qui se répètent beaucoup, mais qui n'ont pas de sens ajouté (Ex. \expword{
	Prépositions})
	\item Les mots qui apparient dans plusieurs documents sont moins importants
\end{itemize}

\begin{block}{Calcul de Tf-Idf}
	\[
	Tf_d(t) =  |\{t_i \in d / t_i = t\}|
	\hskip2cm 
	Df_D(t) = |\{d \in D / t \in d\}|
	\]
	\[Idf_D(t) = \log_{10} \left( \frac{|\{d \in D\}|}{Df_D(t)} \right)\]
	\[Tf\text{-}Idf_{d, D}(t) = Tf_d(t) * Idf_D(t)\]
\end{block}

\end{frame}


\begin{frame}
\frametitle{Sens et désambigüisation : Représentation vectorielle}
\framesubtitle{Tf-idf : Exemple d'une représentation avec Tf}

\begin{exampleblock}{Exemple d'une collection de phrases}
	\begin{itemize}
		\item S1 : un ordinateur peut vous aider
		\item S2 : il peut vous aider et il veut vous aider
		\item S3 : il veut un ordinateur et un ordinateur pour vous
	\end{itemize}
\end{exampleblock}

\begin{center}
	\begin{tabular}{llllllllll}
	\hline\hline
	& un & ordinateur & peut & vous & aider & il & et & veut & pour \\
	\hline
	S1 & 1 & 1 & 1 & 1 & 1 & 0 & 0 & 0 & 0\\
	S2 & 0 & 0 & 1 & 2 & 2 & 2 & 1 & 1 & 0\\
	S3 & 2 & 2 & 0 & 1 & 0 & 1 & 1 & 1 & 1\\
	\hline\hline
\end{tabular}
\end{center}

\end{frame}


\begin{frame}
\frametitle{Sens et désambigüisation : Représentation vectorielle}
\framesubtitle{Tf-idf : Similarité cosinus}

\begin{minipage}{.68\textwidth}
\begin{itemize}
	\item $a$ et $b$ sont deux documents représentés par deux vecteurs $\overrightarrow{a}$ et $\overrightarrow{b}$ respectivement
	\item La représentation peut être les mots du vocabulaire (Tf ou Tf-Idf) ou d'autres caractéristiques du document
	\item La longueur des vecteurs est de $n$
\end{itemize}
\end{minipage}
\begin{minipage}{.3\textwidth}
\hgraphpage{exp-cos.pdf}
\end{minipage}

\begin{block}{Calcul de similarité cosinus entre deux vecteurs}
	\[
	Cos(\theta) = \frac{\overrightarrow{a} \overrightarrow{b}}{||\overrightarrow{a}||\, ||\overrightarrow{b}||}
	= \frac{\sum_{i=1}^{n} a_i b_i}{\sqrt{\sum_{i=1}^{n} a_i^2} \sqrt{\sum_{i=1}^{n} b_i^2}}
	\]
\end{block}

\end{frame}

\subsection{Mot-Mot}

\begin{frame}
\frametitle{Sens et désambigüisation : Représentation vectorielle}
\framesubtitle{Mot-Mot}

\begin{itemize}
	\item On peut représenter un mot par rapport aux autres mot du vocabulaire en utilisant la co-occurrence
	\item Pour représenter les mots d'un vocabulaire $ V $, on doit utiliser une matrice $|V| \times |V|$
	\item Chaque mot est représenté par $|V|$ mots appelés \keyword{le contexte}
	\item La co-occurrence peut être calculée par rapport aux documents, aux phrases ou des fenêtres auteur du mot
	\item La fenêtre peut être 4 mots avant et 4 mots après
\end{itemize}

\end{frame}

\begin{frame}
\frametitle{Sens et désambigüisation : Représentation vectorielle}
\framesubtitle{Mot-Mot : Exemple avec un fenêtre de 1-1}

\begin{exampleblock}{Exemple d'une collection de phrases}
	\begin{itemize}
		\item S1 : un ordinateur peut vous aider
		\item S2 : il peut vous aider et il veut vous aider
		\item S3 : il veut un ordinateur et un ordinateur pour vous
	\end{itemize}
\end{exampleblock}

\begin{center}
	\scriptsize
	\begin{tabular}{llllllllll}
		\hline\hline
		& un & ordinateur & peut & vous & aider & il & et & veut & pour \\
		\hline
		un & 0 & 3 & 0 & 0 & 0 & 0 & 1 & 1 & 0\\
		ordinateur & 3 & 0 & 1 & 0 & 0 & 0 & 1 & 0 & 1\\
		peut & 0 & 1 & 0 & 2 & 0 & 1 & 0 & 0 & 0\\
		vous & 0 & 0 & 2 & 0 & 3 & 0 & 0 & 1 & 1\\
		aider & 0 & 0 & 0 & 3 & 0 & 0 & 1 & 0 & 0\\
		il & 0 & 0 & 1 & 0 & 0 & 0 & 1 & 2 & 0\\
		et & 0 & 0 & 0 & 0 & 1 & 1 & 0 & 0 & 0\\
		veut & 1 & 0 & 0 & 1 & 0 & 2 & 0 & 0 & 0\\
		pour & 0 & 1 & 0 & 1 & 0 & 0 & 0 & 0 & 0\\
		\hline\hline
	\end{tabular}
\end{center}

\end{frame}

\begin{frame}
\frametitle{Sens et désambigüisation : Représentation vectorielle}
\framesubtitle{Mot-Mot : Similarité cosinus}

\begin{itemize}
	\item On peut calculer la similarité entre deux mots
	\item Une mesure de similarité est cosinus (vue précédemment)
\end{itemize}

\hgraphpage{exp-word-v1_.pdf}

\begin{minipage}{.58\textwidth}
	\begin{figure}
		\caption{Un exemple des vecteurs de co-occurrence à partir de Wikipedia et une visualisation de deux mots \cite{2019-jurafsky-martin} }
	\end{figure}
\end{minipage}
\begin{minipage}{.4\textwidth}
	\hgraphpage{exp-word-v2_.pdf}
\end{minipage}

\end{frame}

\subsection{Mot-document}

\begin{frame}
\frametitle{Sens et désambigüisation : Représentation vectorielle}
\framesubtitle{Mot-document}

\begin{itemize}
	\item LSA
\end{itemize}

\end{frame}



%===================================================================================
\section{Word embedding}
%===================================================================================

\begin{frame}
\frametitle{Sens des mots et désambigüisation lexicale}
\framesubtitle{Word embedding}

\begin{itemize}
	\item 	
\end{itemize}

\end{frame}


\subsection{Word2vec}

\begin{frame}
\frametitle{Sens des mots et désambigüisation lexicale : Word embedding}
\framesubtitle{Word2vec}

\begin{itemize}
	\item 
\end{itemize}

\end{frame}

\subsection{GloVe}

\begin{frame}
\frametitle{Sens des mots et désambigüisation lexicale : Word embedding}
\framesubtitle{GloVe}

\begin{itemize}
	\item 
\end{itemize}

\end{frame}

\subsection{BERT}

\begin{frame}
\frametitle{Sens des mots et désambigüisation lexicale : Word embedding}
\framesubtitle{BERT}

\begin{itemize}
	\item 
\end{itemize}

\end{frame}

\subsection{ELMo}

\begin{frame}
\frametitle{Sens des mots et désambigüisation lexicale : Word embedding}
\framesubtitle{ELMo}

\begin{itemize}
	\item 
\end{itemize}

\end{frame}

%===================================================================================
\section{Bases de donnée lexicales}
%===================================================================================

\begin{frame}
\frametitle{Sens des mots et désambigüisation lexicale}
\framesubtitle{Bases de donnée lexicales}

\begin{itemize}
	\item 	
\end{itemize}

\end{frame}

\subsection{Relations sémantiques}

\begin{frame}
\frametitle{Sens des mots et désambigüisation lexicale : BD lexicales}
\framesubtitle{Relations sémantiques}

\begin{itemize}
	\item 
\end{itemize}

\end{frame}

\subsection{Wordnet}

\begin{frame}
\frametitle{Sens des mots et désambigüisation lexicale : BD lexicales}
\framesubtitle{Wordnet}

\begin{itemize}
	\item 
\end{itemize}

\end{frame}

\subsection{Autres ressources}

\begin{frame}
\frametitle{Sens des mots et désambigüisation lexicale : BD lexicales}
\framesubtitle{Autres ressources}

\begin{itemize}
	\item BabelNet, FrameNet
\end{itemize}

\end{frame}

%===================================================================================
\section{Désambigüisation lexicale}
%===================================================================================

\begin{frame}
\frametitle{Sens des mots et désambigüisation lexicale}
\framesubtitle{Désambigüisation lexicale}

\begin{itemize}
	\item 	
\end{itemize}

\end{frame}

\subsection{Algorithme de Lesk}

\begin{frame}
\frametitle{Sens des mots et désambigüisation lexicale : Désambigüisation}
\framesubtitle{Algorithme de Lesk}

\begin{itemize}
	\item 
\end{itemize}

\end{frame}

\subsection{Contextual Embeddings}

\begin{frame}
\frametitle{Sens des mots et désambigüisation lexicale : Désambigüisation}
\framesubtitle{Contextual Embeddings}

\begin{itemize}
	\item 
\end{itemize}

\end{frame}

\insertbibliography{TALN06}{*}

\end{document}

